\chapter[31]{Hyperbolic Geometry and 3-Manifolds}
\label{ch:chow-iv-hyperbolic-geometry-3-manifolds}
\dcref{ch31}
\source{chow-et-al-ricci-flow-part-iv:page-0192}

\section{Introduction to hyperbolic space}

\begin{definition}[Models of hyperbolic space]
\label{def:chow-iv-hyperbolic-models}
\source{chow-et-al-ricci-flow-part-iv:page-0192,chow-et-al-ricci-flow-part-iv:page-0193}
\dcref{ch31:31.1}
\group{Hyperbolic models}

The disk, upper half-space, and hyperboloid models are
\[
 D^n=\{x\in\mathbb R^n:|x|<1\},\quad
 U^n=\{x\in\mathbb R^n:x_n>0\},\quad
 H^n=\{x\in\mathbb R^{n+1}:\langle x,x\rangle_{n,1}=-1, x_{n+1}>0\},
\]
with metrics \(4|dx|^2/(1-|x|^2)^2\), \(|dx|^2/x_n^2\), and the induced
Lorentz metric, respectively.  The maps
\[
 x\mapsto \frac{x}{1+x_{n+1}},\qquad
 x\mapsto 2\frac{x+e_n}{|x+e_n|^2}-e_n
\]
are isometries between the hyperboloid, disk, and upper half-space models.
\end{definition}

\begin{lemma}[Orientation-preserving isometries of the hyperbolic plane]
\label{lem:chow-iv-hyperbolic-plane-isometries}
\source{chow-et-al-ricci-flow-part-iv:page-0193}
\dcref{ch31:31.2}
\group{Hyperbolic isometries}

Every orientation-preserving isometry of \(U^2\) is a Mobius transformation
\[
 z\longmapsto \frac{az+b}{cz+d},\qquad a,b,c,d\in\mathbb R,\quad ad-bc=1.
\]
\end{lemma}

\begin{example}[Vertical geodesic and geodesics in \(U^2\)]
\label{ex:chow-iv-hyperbolic-plane-geodesics}
\source{chow-et-al-ricci-flow-part-iv:page-0194}
\dcref{ch31:31.3}
\group{Hyperbolic isometries}

The positive imaginary axis is minimizing, with
\(d_{U^2}(ia,ib)=\log(b/a)\).  Mobius invariance then shows that all
geodesics in \(U^2\) are vertical lines or circles perpendicular to the real
axis.
\end{example}

\begin{lemma}[Inversion preserves the upper half-space metric]
\label{lem:chow-iv-inversion-hyperbolic-isometry}
\source{chow-et-al-ricci-flow-part-iv:page-0194}
\dcref{ch31:31.4}
\group{Hyperbolic isometries}

The inversion \(x\mapsto x/|x|^2\) preserves \(g_U=|dx|^2/x_n^2\) on \(U^n\).
\end{lemma}

\begin{lemma}[Isometry group of hyperbolic space]
\label{lem:chow-iv-isometry-group-hyperbolic-space}
\source{chow-et-al-ricci-flow-part-iv:page-0195}
\dcref{ch31:31.5}
\group{Hyperbolic isometries}

\(\operatorname{Isom}(\mathbb H^n)\cong\operatorname{Mob}(n-1)\), where every
element is a composition of a dilation, an orthogonal map, a translation, and
possibly an inversion.  The orientation-preserving subgroup is characterized
by the parity of the inversion and the determinant of the orthogonal map.
\end{lemma}

\begin{lemma}[Geodesics in the upper half-space model]
\label{lem:chow-iv-upper-half-space-geodesics}
\source{chow-et-al-ricci-flow-part-iv:page-0196}
\dcref{ch31:31.6}
\group{Hyperbolic isometries}

All geodesics in \(U^n\) are portions of vertical lines or circles meeting the
hyperplane at infinity orthogonally.
\end{lemma}

\begin{definition}[Elliptic, parabolic, and hyperbolic isometries]
\label{def:chow-iv-isometry-types}
\source{chow-et-al-ricci-flow-part-iv:page-0196}
\dcref{ch31:31.7}
\group{Hyperbolic isometries}

An orientation-preserving isometry is elliptic if it fixes a point of
\(\mathbb H^n\), parabolic if it fixes exactly one point at infinity and no
interior point, and hyperbolic if it fixes exactly two points at infinity and
no interior point.
\end{definition}

\begin{example}[Standard isometries]
\label{ex:chow-iv-standard-isometry-types}
\source{chow-et-al-ricci-flow-part-iv:page-0196,chow-et-al-ricci-flow-part-iv:page-0197}
\dcref{ch31:31.8}
\group{Hyperbolic isometries}

Rotations of the disk are elliptic; translations parallel to the boundary in
the upper half-space model are parabolic; and maps
\((x',x_n)\mapsto(\lambda A x',\lambda x_n)\), \(\lambda\ne1\), are hyperbolic.
\end{example}

\begin{lemma}[Trace criterion in dimension three]
\label{lem:chow-iv-trace-isometry-classification}
\source{chow-et-al-ricci-flow-part-iv:page-0197}
\dcref{ch31:31.9}
\group{Hyperbolic isometries}

For \([M]\in\operatorname{PSL}(2,\mathbb C)\), represented by
\(M\in\operatorname{SL}(2,\mathbb C)\setminus\{\pm I\}\), the class is
elliptic iff \(\operatorname{tr}M\in(-2,2)\), parabolic iff
\(\operatorname{tr}M=\pm2\), and hyperbolic iff the trace is real with
absolute value greater than two or is nonreal.  An element fixing infinity
acts by \((z,t)\mapsto(a^2z+ab,|a|^2t)\).
\end{lemma}

\begin{definition}[Nilpotent group]
\label{def:chow-iv-nilpotent-group}
\source{chow-et-al-ricci-flow-part-iv:page-0197}
\dcref{ch31:31.10}
\group{Discrete groups}

For \(G^0=G\) and \(G^k=[G,G^{k-1}]\), a group is nilpotent when
\(G^\ell=\{e\}\) for some \(\ell\).
\end{definition}

\begin{lemma}[Nilpotent subgroups of \(\operatorname{Isom}(\mathbb H^n)\)]
\label{lem:chow-iv-nilpotent-isometry-subgroups}
\source{chow-et-al-ricci-flow-part-iv:page-0198}
\uses{def:chow-iv-isometry-types, def:chow-iv-nilpotent-group}
\dcref{ch31:31.11}
\group{Discrete groups}


If \(G\subset\operatorname{Isom}(\mathbb H^n)\) is nilpotent, then either it is
trivial, contains a nontrivial elliptic element, all nontrivial elements are
hyperbolic with the same two ideal fixed points, or all nontrivial elements are
parabolic with one common ideal fixed point.
\end{lemma}

\begin{definition}[Horoball and horosphere]
\label{def:chow-iv-horoball}
\source{chow-et-al-ricci-flow-part-iv:page-0198}
\dcref{ch31:31.12}
\group{Hyperbolic cusps}

A horoball is the intersection of the disk model with a Euclidean ball tangent
to its ideal boundary; its boundary is a horosphere.  In the upper half-space
model horoballs are also regions \(\{x_n\ge c\}\).
\end{definition}

\section{Topology and geometry of hyperbolic 3-manifolds}

\begin{lemma}[Complete hyperbolic 3-manifolds are irreducible]
\label{lem:chow-iv-hyperbolic-irreducible}
\source{chow-et-al-ricci-flow-part-iv:page-0199,chow-et-al-ricci-flow-part-iv:page-0200}
\dcref{ch31:31.13}
\group{3-manifold topology}

Every embedded two-sphere in a complete hyperbolic three-manifold bounds an
embedded three-ball.
\end{lemma}

\begin{theorem}[Loop theorem]
\label{thm:chow-iv-loop-theorem}
\source{chow-et-al-ricci-flow-part-iv:page-0200}
\dcref{ch31:31.14}
\group{3-manifold topology}

Let \(\Sigma\subset\partial M\) be connected and let \(N\lhd\pi_1(\Sigma)\).
If \(\ker(\pi_1(\Sigma)\to\pi_1(M))\setminus N\ne\varnothing\), there is an
embedded disk whose boundary lies in \(\Sigma\) and represents an element
outside \(N\).
\end{theorem}

\begin{corollary}[Compressible boundary components]
\label{cor:chow-iv-compressible-boundary}
\source{chow-et-al-ricci-flow-part-iv:page-0200}
\uses{thm:chow-iv-loop-theorem}
\dcref{ch31:31.15}
\group{3-manifold topology}


A closed two-sided compressible boundary surface of positive genus contains a
simple closed curve representing a nontrivial element of the kernel of the
inclusion on fundamental groups; this element is primitive.
\end{corollary}

\begin{lemma}[Simple curves represent primitive elements]
\label{lem:chow-iv-simple-curve-primitive}
\source{chow-et-al-ricci-flow-part-iv:page-0200,chow-et-al-ricci-flow-part-iv:page-0201}
\dcref{ch31:31.16}
\group{3-manifold topology}

Every nontrivial element represented by a simple closed curve in an oriented
closed surface is primitive in its fundamental group.
\end{lemma}

\begin{remark}[Hyperbolic proof of primitivity]
\label{rem:chow-iv-geodesic-primitivity}
\source{chow-et-al-ricci-flow-part-iv:page-0201}
\dcref{ch31:31.17}
\group{3-manifold topology}

On a closed hyperbolic surface uniqueness of geodesic representatives gives a
second proof: a simple geodesic cannot be a nontrivial power.
\end{remark}

\begin{lemma}[Kernel for a compressible torus]
\label{lem:chow-iv-compressible-torus-kernel}
\source{chow-et-al-ricci-flow-part-iv:page-0201,chow-et-al-ricci-flow-part-iv:page-0202}
\dcref{ch31:31.18}
\group{3-manifold topology}

If \(T^2\hookrightarrow M^3\) is compressible and \(\pi_1(M)\) is torsion-free,
then the kernel of \(\pi_1(T)\to\pi_1(M)\) is either all of \(\mathbb Z^2\) or
an infinite cyclic subgroup generated by a primitive element.
\end{lemma}

\begin{remark}[Lens spaces]
\label{rem:chow-iv-compressible-torus-lens-space}
\source{chow-et-al-ricci-flow-part-iv:page-0201}
\dcref{ch31:31.19}
\group{3-manifold topology}

The preceding kernel statement fails for lens spaces.
\end{remark}

\begin{lemma}[Boundary inclusion on first homology]
\label{lem:chow-iv-boundary-h1-rank}
\source{chow-et-al-ricci-flow-part-iv:page-0202}
\dcref{ch31:31.20}
\group{3-manifold topology}

If \(M^3\) is compact orientable with boundary \(\Sigma\), then the rank of
\(H_1(\Sigma;\mathbb R)\to H_1(M;\mathbb R)\) equals one half of
\(\dim H_1(\Sigma;\mathbb R)\), by Lefschetz duality and exactness.
\end{lemma}

\begin{lemma}[Separating incompressible surfaces]
\label{lem:chow-iv-separating-incompressible-surfaces}
\source{chow-et-al-ricci-flow-part-iv:page-0203}
\dcref{ch31:31.21}
\group{3-manifold topology}

If a closed surface separates a closed 3-manifold into two pieces and is
incompressible in both, then it is incompressible in the whole manifold and the
fundamental groups of both pieces inject into the whole group.
\end{lemma}

\begin{theorem}[Seifert--van Kampen]
\label{thm:chow-iv-seifert-van-kampen}
\source{chow-et-al-ricci-flow-part-iv:page-0203}
\dcref{ch31:31.22}
\group{3-manifold topology}

For a union \(X=X_1\cup X_2\) with path-connected intersection,
\(\pi_1(X)\) is the free product of \(\pi_1(X_1)\) and \(\pi_1(X_2)\) modulo
the normal subgroup identifying the two inclusion images.  If both maps from
the intersection are injective, the maps from the pieces are injective.
\end{theorem}

\begin{lemma}[Boundary tori are incompressible]
\label{lem:chow-iv-boundary-torus-incompressible}
\source{chow-et-al-ricci-flow-part-iv:page-0203,chow-et-al-ricci-flow-part-iv:page-0204}
\dcref{ch31:31.25}
\group{Hyperbolic 3-manifolds}

Every torus slice in a finite-volume hyperbolic cusp is incompressible in the
ambient hyperbolic 3-manifold.
\end{lemma}

\begin{lemma}[Embedded tori in finite-volume hyperbolic 3-manifolds]
\label{lem:chow-iv-embedded-torus-trichotomy}
\source{chow-et-al-ricci-flow-part-iv:page-0204,chow-et-al-ricci-flow-part-iv:page-0205,chow-et-al-ricci-flow-part-iv:page-0206}
\dcref{ch31:31.27}
\group{Hyperbolic 3-manifolds}

An embedded torus is either incompressible and isotopic to a standard cusp
slice, compressible and bounds a solid torus, or compressible and bounds a
compact 3-manifold contained in an embedded 3-ball.
\end{lemma}

\begin{example}[Compressible torus not bounding a solid torus]
\label{ex:chow-iv-knotted-compressible-torus}
\source{chow-et-al-ricci-flow-part-iv:page-0206}
\dcref{ch31:31.28}
\group{Hyperbolic 3-manifolds}

A knotted torus in a ball embedded in a hyperbolic 3-manifold gives an example
of the third alternative in Lemma~31.27.
\end{example}

\begin{definition}[Topological hyperbolic piece]
\label{def:chow-iv-topological-hyperbolic-piece}
\source{chow-et-al-ricci-flow-part-iv:page-0206}
\dcref{ch31:31.29}
\group{Hyperbolic 3-manifolds}

A compact submanifold of a closed 3-manifold whose boundary is a union of tori
and whose interior carries a complete finite-volume hyperbolic metric is a
topological hyperbolic piece.
\end{definition}

\begin{corollary}[Incompressibility of hyperbolic pieces]
\label{cor:chow-iv-hyperbolic-piece-incompressibility}
\source{chow-et-al-ricci-flow-part-iv:page-0206}
\uses{def:chow-iv-topological-hyperbolic-piece, lem:chow-iv-separating-incompressible-surfaces, lem:chow-iv-boundary-torus-incompressible}
\dcref{ch31:31.30}
\group{Hyperbolic 3-manifolds}


If the boundary of a topological hyperbolic piece is incompressible in its
complement, then it is incompressible in the closed ambient manifold.
\end{corollary}

\section{The Margulis lemma and hyperbolic cusps}

\begin{definition}[Topological end]
\label{def:chow-iv-topological-end}
\source{chow-et-al-ricci-flow-part-iv:page-0206,chow-et-al-ricci-flow-part-iv:page-0207}
\dcref{ch31:31.31}
\group{Hyperbolic ends}

A compact set \(K\) is end-complementary if components of complements stabilize
under enlargement of \(K\).  Such stabilized components are the topological
ends of a noncompact manifold.
\end{definition}

\begin{example}[Cylinder]
\label{ex:chow-iv-cylinder-ends}
\source{chow-et-al-ricci-flow-part-iv:page-0207}
\dcref{ch31:31.32}
\group{Hyperbolic ends}

\(N^{n-1}\times\mathbb R\) has two topological ends.
\end{example}

\begin{definition}[Hyperbolic cusp]
\label{def:chow-iv-hyperbolic-cusp}
\source{chow-et-al-ricci-flow-part-iv:page-0207}
\dcref{ch31:31.34}
\group{Hyperbolic cusps}

For a flat \((n-1)\)-manifold \((\mathcal V,g_{\rm flat})\), a hyperbolic cusp is
\[
([0,\infty)\times\mathcal V,\ g_{\rm cusp}=dr^2+e^{-2r}g_{\rm flat}).
\]
\end{definition}

\begin{theorem}[Margulis lemma, algebraic form]
\label{thm:chow-iv-margulis-algebraic}
\source{chow-et-al-ricci-flow-part-iv:page-0208}
\dcref{ch31:31.36}
\group{Margulis lemma}

For every \(n\) there is \(\varepsilon_n>0\) such that
\(\{\gamma:d(\gamma x,x)\le\varepsilon_n\}\) is virtually nilpotent, with
index bounded in terms of \(n\) only.
\end{theorem}

\begin{remark}[Three-dimensional Margulis groups]
\label{rem:chow-iv-margulis-three-dimensional}
\source{chow-et-al-ricci-flow-part-iv:page-0208}
\dcref{ch31:31.37}
\group{Margulis lemma}

For \(n=3\), the Margulis group is virtually abelian.
\end{remark}

\begin{theorem}[Margulis lemma, local form]
\label{thm:chow-iv-margulis-local}
\source{chow-et-al-ricci-flow-part-iv:page-0208}
\dcref{ch31:31.38}
\group{Margulis lemma}

For \(0<\varepsilon\le\varepsilon_n\), the ball of radius \(\varepsilon/2\) is
isometric to the corresponding quotient by the subgroup generated by loops of
length at most \(\varepsilon\); that subgroup is trivial, infinite cyclic, or
a discrete free Euclidean isometry group.
\end{theorem}

\begin{definition}[Thick-thin decomposition]
\label{def:chow-iv-thick-thin-parts}
\source{chow-et-al-ricci-flow-part-iv:page-0209}
\dcref{ch31:31.39}
\group{Margulis lemma}

The \(\varepsilon\)-thin part consists of points supporting a nontrivial loop
of length less than \(\varepsilon\); the thick part consists of points where
all such loops have length at least \(\varepsilon\).
\end{definition}

\begin{example}[Thick and thin parts]
\label{ex:chow-iv-thick-thin-examples}
\source{chow-et-al-ricci-flow-part-iv:page-0209}
\dcref{ch31:31.40}
\group{Margulis lemma}

Simply connected manifolds are entirely thick; closed manifolds are thick for
sufficiently small \(\varepsilon\); and sufficiently far out in a cusp every
point is \(\varepsilon\)-thin.
\end{example}

\begin{lemma}[Injectivity radius and compact thick part]
\label{lem:chow-iv-thick-injectivity-radius}
\source{chow-et-al-ricci-flow-part-iv:page-0209}
\dcref{ch31:31.41}
\group{Margulis lemma}

The injectivity radius on the \(\varepsilon\)-thick part is at least
\(\varepsilon/2\).  In a finite-volume hyperbolic manifold the thick part is
compact.
\end{lemma}

\begin{lemma}[Topological ends are essentially thin]
\label{lem:chow-iv-topological-ends-thin}
\source{chow-et-al-ricci-flow-part-iv:page-0210}
\dcref{ch31:31.42}
\group{Margulis lemma}

In a finite-volume hyperbolic manifold, every topological end contains a
cofinal \(\varepsilon\)-thin end, and for sufficiently small \(\varepsilon\)
every \(\varepsilon\)-thin end is topological.
\end{lemma}

\begin{theorem}[Geometric structure of the thin part]
\label{thm:chow-iv-thin-part-structure}
\source{chow-et-al-ricci-flow-part-iv:page-0210}
\dcref{ch31:31.43}
\group{Margulis lemma}

For \(n\ge3\) and \(\varepsilon\le\varepsilon_n\), the thin part is a disjoint
union of cusp ends, Margulis tubes around short geodesics, and (at exceptional
lengths) geodesic circles.  Cusp pieces have metric
\(dr^2+e^{-2r}g_{\rm flat}\), while tubes are modeled on
\(\sum_i(x^i ds+dx^i)^2+ds^2\).
\end{theorem}

\begin{theorem}[Finite-volume ends are standard cusps]
\label{thm:chow-iv-finite-volume-ends-cusps}
\source{chow-et-al-ricci-flow-part-iv:page-0210,chow-et-al-ricci-flow-part-iv:page-0211}
\uses{lem:chow-iv-topological-ends-thin, thm:chow-iv-thin-part-structure}
\dcref{ch31:31.44}
\group{Hyperbolic cusps}


Every topological end of a finite-volume hyperbolic manifold contains a
cofinal subset isometric to a standard cusp
\([0,\infty)\times\mathcal V\) with metric
\(dr^2+e^{-2r}g_{\rm flat}\), where \(\mathcal V\) is closed and flat.
\end{theorem}

\begin{example}[Margulis tube as a quotient]
\label{ex:chow-iv-margulis-tube-quotient}
\source{chow-et-al-ricci-flow-part-iv:page-0211}
\dcref{ch31:31.45}
\group{Margulis lemma}

The change of variables \(y'=e^s x\), \(y^n=e^s\) identifies the upper
half-space metric with \(\sum_i(x^i ds+dx^i)^2+ds^2\), and a tube is the quotient
of \(B^{n-1}(r)\times\mathbb R\) by \(s\mapsto s+\ell\).
\end{example}

\begin{theorem}[Volume bounds the number of cusp ends]
\label{thm:chow-iv-cusp-volume-bound}
\source{chow-et-al-ricci-flow-part-iv:page-0212}
\dcref{ch31:31.46}
\group{Hyperbolic cusps}

If a finite-volume hyperbolic 3-manifold has \(m\) cusp ends, then
\(\operatorname{Vol}(\mathcal H)\ge m\nu\), where \(\nu\) is the volume of the
regular ideal tetrahedron; moreover it contains \(m\) disjoint cusp ends each
of volume at least \(\sqrt3/4\).
\end{theorem}

\begin{corollary}[Small almost-flat tori lie in cusps]
\label{cor:chow-iv-small-almost-flat-tori}
\source{chow-et-al-ricci-flow-part-iv:page-0212}
\uses{thm:chow-iv-finite-volume-ends-cusps}
\dcref{ch31:31.47}
\group{Hyperbolic cusps}


For every \(C<\infty\), sufficiently small intrinsic sectional curvature and
area, together with \(|\mathrm{II}|\le C\), force an embedded 2-torus in a
finite-volume hyperbolic 3-manifold to lie in a maximal cusp end.
\end{corollary}

\begin{remark}[Use in Ricci-flow limits]
\label{rem:chow-iv-small-tori-application}
\source{chow-et-al-ricci-flow-part-iv:page-0212}
\dcref{ch31:31.48}
\group{Hyperbolic cusps}

Corollary~31.47 is used to locate almost-flat torus cross-sections in noncompact
hyperbolic limits of nonsingular Ricci flows.
\end{remark}

\section{Mostow rigidity}

\begin{theorem}[Mostow rigidity]
\label{thm:chow-iv-mostow-rigidity}
\source{chow-et-al-ricci-flow-part-iv:page-0213}
\dcref{ch31:31.49}
\group{Mostow rigidity}

If two finite-volume hyperbolic manifolds of dimension at least three have
isomorphic fundamental groups, there is a unique isometry inducing the given
isomorphism.
\end{theorem}

\begin{corollary}[Homotopy, diffeomorphism, and isometry]
\label{cor:chow-iv-mostow-equivalences}
\source{chow-et-al-ricci-flow-part-iv:page-0213}
\uses{thm:chow-iv-mostow-rigidity}
\dcref{ch31:31.50}
\group{Mostow rigidity}


For finite-volume hyperbolic manifolds of dimension at least three, being
homotopy equivalent, diffeomorphic, and isometric are equivalent conditions;
volume is therefore a topological invariant.
\end{corollary}

\begin{corollary}[Isometry group]
\label{cor:chow-iv-isometry-group}
\source{chow-et-al-ricci-flow-part-iv:page-0213}
\uses{thm:chow-iv-mostow-rigidity}
\dcref{ch31:31.51}
\group{Mostow rigidity}


The isometry group of a finite-volume hyperbolic manifold of dimension at least
three is finite and isomorphic to \(\operatorname{Out}(\pi_1(\mathcal H))\).
\end{corollary}

\section{Seifert fibered manifolds and graph manifolds}

\begin{definition}[Foliation]
\label{def:chow-iv-foliation}
\source{chow-et-al-ricci-flow-part-iv:page-0214}
\dcref{ch31:31.52}
\group{Collapsing 3-manifolds}

A smooth \(k\)-foliation is an atlas whose transition maps preserve the first
\(n-k\) coordinates; connected maximal unions of plaques are its leaves.
\end{definition}

\begin{definition}[Seifert fibered manifold]
\label{def:chow-iv-seifert-fibered}
\source{chow-et-al-ricci-flow-part-iv:page-0214}
\dcref{ch31:31.53}
\group{Collapsing 3-manifolds}

A Seifert fibered 3-manifold is a compact 3-manifold carrying a foliation by
circles.  Solid tori with the model foliations \(\mathbf F_{p,q}\) are the
building blocks; finite covers are products or circle bundles over surfaces.
\end{definition}

\begin{definition}[Graph manifold]
\label{def:chow-iv-graph-manifold}
\source{chow-et-al-ricci-flow-part-iv:page-0215}
\dcref{ch31:31.54}
\group{Collapsing 3-manifolds}

A graph manifold is obtained from Seifert fibered vertex pieces and
\(T^2\times[0,1]\) edge pieces glued along boundary tori according to a finite
graph.
\end{definition}

\begin{theorem}[Cheeger--Gromov collapsing theorem]
\label{thm:chow-iv-cheeger-gromov-graph-manifold}
\source{chow-et-al-ricci-flow-part-iv:page-0215}
\dcref{ch31:31.55}
\group{Collapsing 3-manifolds}

If a compact 3-manifold with empty or torus boundary admits metrics satisfying
\[
\max_x\operatorname{inj}_{g_i}(x)^2\,\max_M|\operatorname{Rm}_{g_i}|\longrightarrow0,
\]
then it admits an \(F\)-structure and hence is a graph manifold.
\end{theorem}

\begin{claim}[Exercise 31.35: cusp slices]
\label{ex:chow-iv-cusp-slices-umbilic}
\dcref{ch31:31.35}
\group{Hyperbolic cusps}
\source{chow-et-al-ricci-flow-part-iv:page-0208}
\uses{def:chow-iv-hyperbolic-cusp}
The slices \(\{r\}\times\mathcal V\) are totally umbilic and, with normal
\(\partial_r\), satisfy \(\mathrm{II}=-g_{\rm cusp}|_{\{r\}\times\mathcal V}\).
\end{claim}

\begin{claim}[Exercise 31.33: ends under limits]
\label{ex:chow-iv-ends-cheeger-gromov}
\dcref{ch31:31.33}
\group{Hyperbolic ends}
\source{chow-et-al-ricci-flow-part-iv:page-0207}
The number of topological ends need not be lower semicontinuous under pointed
Cheeger--Gromov convergence; manifolds with increasingly many ends can converge
to Euclidean space.
\end{claim}

\begin{claim}[Exercise 31.1: model equivalence]
\label{ex:chow-iv-hyperbolic-model-isometries}
\dcref{ch31:31.1-exercise}
\group{Hyperbolic models}
\source{chow-et-al-ricci-flow-part-iv:page-0193}
\uses{def:chow-iv-hyperbolic-models}
Directly verify that the two displayed maps pull the corresponding metrics
back to one another.
\end{claim}
